\section{Síntesis y ampliación}

\subsection{Mensaje central}

\begin{enumerate}
    \item La capacidad matemática actual de los LLM depende principalmente del SFT y el RL, limitada por el volumen de datos y la posibilidad de verificación
    \item El razonamiento formalizado (Lean) ofrece una señal de verificación perfecta, lo que puede aliviar la escasez de datos y la dificultad de evaluación
    \item LeanDojo ofrece una infraestructura de código abierto para la demostración de teoremas
    \item Los métodos específicos de dominio (como LiPS) pueden superar a los humanos y a los LLM de vanguardia en tareas específicas
    \item La formalización automática (de enunciados y de demostraciones) presenta retos particulares en cada caso
\end{enumerate}

\subsection{Lecturas adicionales}
\begin{itemize}
    \item LeanDojo (NeurIPS 2023): datos y herramientas de código abierto para la demostración de teoremas
    \item LiPS: método híbrido simbólico y neuronal para la demostración de desigualdades
    \item Artículo de postura ``AI for Math'' (arXiv)
    \item Referencia Frontier Math (Epoch AI)
    \item DeepSeek-R1: entrenamiento con RL de un modelo de razonamiento
\end{itemize}

\end{document}
